\section{Problems addressed and mathematical status}
\label{sec:scope}

The current ProveIt manuscript develops polylogarithm identities
together with harmonic-tail, Gamma, Hurwitz, and generalized Stieltjes
calculi \cite{repo}. Recent incoming reports contain both explicit open
questions and extensive proved developments. A continuation must therefore
identify its exact starting point before claiming an advance.

The principal target here is Question 2 in the research section of
\emph{Ordered Boundary Identities and Gamma Kernels} \cite{repoordered}.
That report constructs a Gamma integral basis for a fixed-gap harmonic
recurrence and proves that its initial determinant equals
$(-1)^{m(m-1)/2}K_m$, where $K_m$ is a positive integral. It asks for
closed evaluations of $K_2,K_3,\ldots$. We give one formula for every
$m$, rather than a table of separately evaluated integrals.

The second target is Question 4 of \emph{Balanced Stieltjes Moments}
\cite{repobalanced}: a colored or Lerch analogue of the homogeneous
balancing and its polynomial moment reductions. The nontrivial color
removes the order pole, so the correct generator has no separate
$1/z$ subtraction. This changes the degree of the negative-order
coordinate and makes the limit of color one singular. We prove the
colored theorem on its full natural slit domain and give its exact
connection with the uncolored normalization.

The repeated-mass theorem answers a two-cluster part of the spectral
face question in \cite{repolattice}. It reduces arbitrary multiplicities
to finitely many multiple polylogarithms. The number of lattice terms
remains infinite: the finite dictionary describes the types and depths
of polylogarithms in each term, and does not give a finite evaluation
of the complete lattice sum.

\subsection{A concise result ledger}

\begin{longtable}{@{}>{\raggedright\arraybackslash}p{0.25\textwidth}>{\raggedright\arraybackslash}p{0.45\textwidth}>{\raggedright\arraybackslash}p{0.24\textwidth}@{}}
\toprule
Result & Exact advance & Status\\
\midrule
\endhead
Gamma normalization &
Closed $K_m$ for all $m$; explicit two-step recurrence &
Incoming Question 2 resolved\\[3pt]
Trigamma moments &
After normalization, the first $n+1$ moments are exactly Gaussian
for every positive shift &
Proved consequence of a parameter identity\\[3pt]
Colored balancing &
Unique coefficients in a specified homogeneous ansatz; all integer
moments and exceptional resonances &
Colored counterpart resolved in that ansatz\\[3pt]
Color-one limit &
Explicit pole subtraction and a shift-dependent constant-term
conversion &
Proved normalization identity\\[3pt]
Two mass groups &
Finite multiple-polylogarithm formula at every pair of multiplicities &
Two-cluster sector resolved\\[3pt]
Depth-two parity &
Exact double-shuffle and distribution certificate for one
printed example &
Classical parity made constructive\\[3pt]
$S_6$, revised $S_8$ &
No exact certificate supplied here &
Remain conjectural\\[3pt]
Unequal-tail prefix &
Initial basis normalization is evaluated; later determinant
nonvanishing is a separate issue &
Unrestricted problem remains open\\
\bottomrule
\end{longtable}

\subsection{Conventions and novelty}

The Gamma function and its logarithmic derivative are denoted by
$\Gamma$ and $\psi$; $\psi^{(j)}$ always means argument differentiation.
An order derivative of Hurwitz zeta or Lerch's transcendent is written
$\partial_s^j$. The generalized Stieltjes normalization is
\begin{equation}\label{eq:stieltjes-convention}
 \zeta(1+z,a)=\frac1z+
  \sum_{j=0}^{\infty}\frac{(-1)^j\gamma_j(a)}{j!}z^j,
 \qquad \gamma_0(a)=-\psi(a).
\end{equation}
All real-variable improper integrals asserted as ordinary integrals
converge absolutely. A constant term is always taken in a specified
parameter; no unspecified regularization is used.

Gustafson's integral, the Gamma reflection formula, the Lerch integral
representation, Bernoulli and Euler special values, and
multiple-polylogarithm Landen transformations are classical
\cite{gustafson1994,derkachov2017,dlmfgamma,dlmflerch,dlmfhurwitz,okudaueno}.
Our normalization theorem is a new application to the specific
incoming basis; the colored and clustered statements are explicit
extensions of the inspected incoming formulas. The finite trigamma
moment law is derived here from the complete parameter identity.
A targeted literature search supports these comparisons but does not
establish worldwide historical priority. No arithmetic independence,
minimal period basis, or proof-assistant formalization is claimed.

